\documentclass[a4paper,12pt]{extarticle}
\usepackage{graphicx}

% Пакеты для русского языка и кодировки
\usepackage[T2A]{fontenc}
\usepackage[utf8]{inputenc}
\usepackage[russian]{babel}

% Пакеты для математики и теорем
\usepackage{amsmath, amssymb, amsthm}

% Пакеты для графики и гиперссылок
\usepackage{hyperref}

% Настройка полей
\usepackage{geometry}
\geometry{top=2cm, bottom=2cm, left=2.5cm, right=1.5cm}

% Для оформления кода и таблиц
\usepackage{listings}
\usepackage{longtable,booktabs,calc}
\usepackage[table]{xcolor}

% Настройка межстрочного интервала
\usepackage{setspace}
\onehalfspacing

% Отключает нумерацию для part, chapter, section и т.д.
\setcounter{secnumdepth}{-3}

\hypersetup{
    colorlinks=true,
    linkcolor=black,
    urlcolor=blue
}

\usepackage{booktabs}
\geometry{top=18mm,bottom=18mm,left=20mm,right=20mm}
\setstretch{1.05}

\begin{document}

\begin{center}
    {\Large\bfseries Ручной перевод $255.9999997_{10}$ в IEEE 754 \texttt{float32}}\\[3mm]
    Округление к ближайшему; при равенстве расстояний --- к чётному
\end{center}

\section*{1. Поля формата и знак}
Число \texttt{float32} занимает 32 бита: 1 бит знака $S$, 8 битов
смещённого порядка $e$ и 23 бита дробной части значащей части $M$.
Для нормализованного числа значение равно
\[
    (-1)^S\,(1.M)_2\,2^{e-127}.
\]
Исходное число положительно, поэтому $\boxed{S=0}$.

\section*{2. Перевод в двоичную систему}
Целая часть:
\[
    255=128+64+32+16+8+4+2+1=11111111_2.
\]
Для дробной части $0.9999997$ умножаем остаток на 2; целая часть
произведения даёт следующий двоичный разряд:
\begin{center}
\begin{tabular}{@{}rllc@{}}
\toprule
Шаг & Остаток до умножения & Удвоенный остаток & Бит \\
\midrule
1  & $0.9999997$ & $1.9999994$ & 1 \\
2  & $0.9999994$ & $1.9999988$ & 1 \\
3  & $0.9999988$ & $1.9999976$ & 1 \\
4  & $0.9999976$ & $1.9999952$ & 1 \\
$\cdots$ & $\cdots$ & $\cdots$ & $\cdots$ \\
16 & $0.9901696$ & $1.9803392$ & 1 \\
17 & $0.9803392$ & $1.9606784$ & 1 \\
18 & $0.9606784$ & $1.9213568$ & 1 \\
\bottomrule
\end{tabular}
\end{center}
Почему пропущенные разряды тоже единицы? Разность до~1 равна
$0.0000003 < 2^{-18}$. Следовательно, первые 18 двоичных цифр дробной
части --- единицы: при меньшей дроби среди первых 18 цифр появился бы ноль.

\clearpage
\section*{3. Нормализация и 23 хранимых бита}
Передвигаем точку в записи $11111111.\ldots_2$ на семь позиций влево:
\[
    11111111.\ldots_2=1.1111111\ldots_2\cdot 2^7.
\]
После скрытой ведущей единицы в поле $M$ попадают 7 единиц из целой
части и следующие 16 единиц из дробной: $7+16=23$.
\[
    \mathtt{1.}\underbrace{\mathtt{11111111111111111111111}}_{23\ \text{хранимых бита}}
    \mathtt{11}\ldots\times 2^7
\]
Первый отбрасываемый бит (17-й из дробной части) равен 1, следующий
тоже равен 1. Значит, отбрасываемая часть больше половины шага:
округляем вверх. До округления смещённый порядок был бы
$7+127=134=10000110_2$.

\section*{4. Перенос при округлении}
Прибавление единицы к последнему хранимому биту переносится через все
23 единицы:
\[
\begin{array}{r@{}l}
    1.&\texttt{11111111111111111111111}\\
 +\;0.&\texttt{00000000000000000000001}\\ \hline
   10.&\texttt{00000000000000000000000}
\end{array}
\]
Получается $10.0\ldots_2\cdot2^7=1.0\ldots_2\cdot2^8=256_{10}$.
Порядок увеличивается на единицу:
\[
    E=8,\qquad e=E+127=135=\texttt{10000111}_2,
    \qquad M=\texttt{00000000000000000000000}.
\]

\section*{5. Итоговое слово и проверка}
Объединяем знак, порядок и дробное поле:
\begin{center}
    \texttt{0 | 10000111 | 00000000000000000000000}\\[2mm]
    \texttt{0100 0011 1000 0000 0000 0000 0000 0000}
\end{center}
По четыре бита на шестнадцатеричную цифру дают
\[
    \boxed{\texttt{43800000}_{16}}.
\]
Непосредственная проверка округления: ближайшее меньшее значение
\texttt{float32} равно
\[
    256-2^{-16}=255.9999847412109375.
\]
Расстояние до него составляет $0.0000149587890625$, а расстояние до
$256$ составляет всего $0.0000003$. Поэтому ответом действительно служит
\[
    \boxed{\operatorname{float32}(255.9999997)=256}.
\]

\end{document}
